\documentclass[11pt,a4paper,twocolumn]{article}

% === Core Packages ===
\usepackage[margin=2cm]{geometry}
\usepackage[utf8]{inputenc}
\usepackage[T1]{fontenc}
\usepackage{times}
\usepackage{amsmath,amssymb}
\usepackage{graphicx}
\usepackage{booktabs}
\usepackage{float}
\usepackage{enumitem}
\usepackage{tikz}
\usepackage{pgfplots}
\pgfplotsset{compat=1.18}
\usetikzlibrary{arrows.meta,shapes.geometric,positioning,calc,decorations.markings,fit,backgrounds}

% === Citation (IEEE style) ===
\usepackage[numbers,sort&compress]{natbib}
\bibliographystyle{IEEEtranN}

% === Hyperref ===
\usepackage[hidelinks]{hyperref}
\usepackage{xurl}

% === Settings ===
\setlength{\parindent}{1em}
\setlength{\parskip}{0.3em}

% === PDF metadata ===
\hypersetup{
  pdfauthor={Pavel Kudrna},
  pdfsubject={Personix — a decentralized reputation network},
  pdfkeywords={Personix, decentralized reputation, reputation social network,
    decentralized identifiers, DID, meritocracy, censorship resistance},
}

% === Title ===
\title{\textbf{شبكة السمعة اللامركزية:\\إطار عمل للتنظيم المجتمعي الطبيعي}}
\author{Pavel Kudrna\\Personix Project --- \url{https://personix.org}\\Prague, Czechia\\المسودة الأولى --- مارس 2026}
\date{}

\begin{document}
\maketitle

% ============================================================
% ABSTRACT
% ============================================================
\begin{abstract}
إن الإحساس بالعدالة---أي القناعة بأن الأخطاء تُصحَّح وأن الجدارة تُكافَأ---هو أقوى قوة تماسك في المجتمع. وعندما تعجز مؤسسات الحوكمة المركزية عن تحقيق هذا الإحساس لأن تكلفة إنفاذ حقٍّ ما تفوق قيمة الحق نفسه، يتآكل التماسك الاجتماعي. تفشل الهياكل المؤسسية الحالية في حل فئة كبيرة من النزاعات بالطريقة المنهجية نفسها التي يفشل بها التخطيط المركزي في تخصيص الموارد: من خلال عدم تماثل المعلومات، وغياب حلقات التغذية الراجعة، وانفصال صناع القرار عن تبعات قراراتهم. تقدّم هذه الورقة شبكة السمعة الاجتماعية (RSN): طبقة اتصال لامركزية، غير قابلة للفساد، ومقاومة للرقابة، مبنية على المعرّفات اللامركزية (DIDs) التي تُمكّن المشاركين المستعارين من نشر معلومات السمعة المتعلقة بالسلوك في العالم الحقيقي والتحقق منها واستهلاكها. تستند الآلية الأساسية إلى ثلاثة مبادئ تصميمية: (1)~بنية تكلفة غير متماثلة يكون فيها قراءة بيانات السمعة رخيصًا بينما تتطلب الكتابة جهدًا قابلًا للتحقق؛ (2)~بروتوكول اختيار مُحقِّق غير حتمي يقوم على التجزئة المتسقة ويُسعّر الادعاءات الجذرية عبر تكاليف تكرار متصاعدة؛ و(3)~نموذج سلطة سمعة مدفوع بالسوق يراهن فيه المحقّقون الخاصون بسمعتهم المتراكمة على دقة الادعاءات التي يصادقون عليها. تُنتج المعمارية الناتجة جدارةً استحقاقية ناشئة دون تنسيق مركزي، وتُتيح مسار انتقال تدريجيًا وقابلًا للعكس من الأنظمة القائمة التي تديرها الدولة.
\end{abstract}

% ============================================================
% 1. INTRODUCTION
% ============================================================
\section{المقدمة}

\subsection{مشكلة الثقة المركزية}

تقوم الحوكمة الحديثة على افتراض أساسي: أن المؤسسات المركزية قادرة على التوسّط في الثقة بين الغرباء على نطاق واسع. فالمحاكم تفصل في النزاعات. والسجلات تُصادق على الملكية. والهيئات التنظيمية تُنفّذ المعايير. لقد صُممت هذه المؤسسات في عصر كانت فيه المعلومات نادرة، والاتصال بطيئًا، وكان تفويض السلطة إلى هيئة مركزية هو آلية التنسيق الممكنة الوحيدة. غير أن الحوكمة المركزية تفشل للأسباب البنيوية نفسها التي يفشل بها التخطيط المركزي: عدم تماثل المعلومات، وغياب حلقات التغذية الراجعة، وانفصال صناع القرار عن تبعات قراراتهم.

يفشل هذا الافتراض، على سبيل المثال، عندما تفوق تكلفة الوساطة المؤسسية قيمة النزاع. لنتأمل مستأجرًا يكتشف أن شقته المستأجرة تفتقر إلى شهادة سلامة كهربائية مطلوبة قانونًا---وهو ظرف يشكّل خطرًا مباشرًا على الحياة. إن حق المستأجر القانوني في الانسحاب من العقد لا لبس فيه. ومع ذلك فإن تكلفة إنفاذ هذا الحق عبر النظام القضائي تفوق القيمة المالية للوديعة والأضرار المستحقة، حتى مع احتساب التعويض القانوني المحتمل~\cite{czcourtfees}. والاستجابة العقلانية هي تحمّل الخسارة والانسحاب.

هذه ليست حالة استثنائية مرَضية. إنها التجربة الافتراضية لفئة كبيرة من النزاعات يكون فيها الضرر حقيقيًا لكن المخاطر المالية تقع دون العتبة التي يصبح عندها الإنفاذ المؤسسي عقلانيًا اقتصاديًا. والنتيجة نظامٌ يُحابي بنيويًا الفاعلين السيئين: فمن يؤذون الآخرين بأحجام دون هذه العتبة يمكنهم التصرف بلا عقاب، لأن تكلفة السعي إلى العدالة تقع على الطرف المتضرر.

المثال أعلاه مأخوذ من البيئة القانونية للمؤلف---النظام المدني القانوني التشيكي. وفي التقاليد القانونية الأخرى---القانون العام القائم على السوابق، والقانون العرفي، والأنظمة المختلطة---تفشل العدالة بطرق مختلفة وبدرجات متفاوتة، تبعًا للخصوصيات الثقافية والإجرائية للولاية القضائية. ومن المرجح أن يتعرّف القرّاء على أمثلة مماثلة من سياقاتهم الخاصة. أما ما هو عالمي بنيويًا فهو النمط: النزاعات التي تقع قيمتها دون عتبة الإنفاذ العقلاني اقتصاديًا تنتهي بتحمّل الطرف المتضرر للخسارة. لا تَعِد شبكة RSN بعدالة كاملة---إنها فقط تقلّص الميزة البنيوية التي يتمتع بها الفاعلون السيئون عندما يكون الإنفاذ المؤسسي غير اقتصادي.

\subsection{لماذا تقصّر الحلول القائمة}

\textbf{أطر الهوية اللامركزية (DID)}~\cite{did2022} توفّر آليات تشفيرية لإدارة الهوية ذاتية السيادة لكنها تركّز بالأساس على التحقق من الهوية دون معالجة المسألة الأوسع المتعلقة بالسمعة السلوكية.

\textbf{الرموز المرتبطة بالروح (SBTs)}~\cite{weyl2022} تلتقط الفكرة القائلة إن السمعة غير قابلة للتحويل لكنها تعتمد على البنية التحتية لسلسلة الكتل مع قيود على قابلية التوسّع، ولا تعالج الحوافز الاقتصادية التي تحكم إدخال المعلومات. أما المفاهيم ذات الصلة مثل رموز الجدارة (كما في ليبرلاند) فتتشارك فلسفة مماثلة لكنها تظل مقيّدة بأطر ولاية قضائية محددة.

\textbf{المنظمات المستقلة اللامركزية (DAOs)}~\cite{buterin2014} تُطبّق الحوكمة عبر العقود الذكية لكنها تُعيد إنتاج ديناميكيات بلوتوقراطية يكون فيها النفوذ متناسبًا مع رأس المال. والتصويت التربيعي~\cite{buterin2019} يخفّف ذلك جزئيًا لكنه يحدّ من التبنّي العملي. كما تواجه المنظمات المستقلة اللامركزية \emph{مشكلة الأوراكل} الأساسية: العقود الذكية لا تستطيع التحقق بموثوقية من حالات العالم الحقيقي دون مصدر خارجي موثوق، ولهذه المشكلة لا يوجد حل عام ضمن معمارية سلسلة الكتل.

لا يجمع أي نظام قائم بين اللامركزية، ومقاومة الرقابة، والاستعارة، وقابلية التحقق من تكلفة الدخول، والإجماع الناشئ.

\subsection{المساهمات}

تقدّم هذه الورقة المساهمات التالية:
\begin{enumerate}[nosep]
\item تعريف \textbf{شبكة السمعة الاجتماعية (RSN)}، وهي بروتوكول لامركزي يوسّع إطار DID لدعم تبادل السمعة الثنائي.
\item \textbf{بروتوكول اختيار مُحقِّق غير حتمي} قائم على التجزئة المتسقة~\cite{karger1997}.
\item \textbf{مبدأ التطرف المكلف}، الذي يُسعّر الادعاءات وفقًا لبُعدها عن إجماع الشبكة عبر ثلاث قنوات تكلفة متراكمة.
\item \textbf{نموذج السلطة ذات السمعة} بحوافز غير متماثلة تكون فيها تكلفة المصادقة الكاذبة أعلى كارثيًا من الرسوم المشروعة.
\item \textbf{مسار انتقال تدريجي} عبر بنية مالية موازية.
\end{enumerate}

% ============================================================
% 2. DESIGN PRINCIPLES
% ============================================================
\section{مبادئ التصميم}

تُقيَّد معمارية RSN بخمسة مبادئ تصميم غير قابلة للتفاوض.

\textbf{الاستمرارية والتطوّر.} يتعايش النظام مع المؤسسات القائمة خلال فترة انتقالية غير محددة المدة. ويجب أن يكون الانتقال قابلًا للعكس في كل مرحلة.

\textbf{الطوعية والمسؤولية.} المشاركة طوعية، لكن الطوعية مقترنة بالمسؤولية: فالمشارك الذي يتولى السيطرة على وظيفة مؤسسية يتحمّل في الوقت نفسه الالتزامات المصاحبة لها.

\textbf{التنوع والحياد الثقافي.} ينشأ الإجماع من التفاعلات الثنائية، لا من قواعد مسبقة يفرضها مصمّمو النظام.

\textbf{المرونة ومقاومة الرقابة.} يجب أن تفوق تكلفة الرقابة على الشبكة تكلفة تحمّلها. فقط الشمولية المطلقة---بتكلفة سياسية واقتصادية مدمّرة---يمكنها قمع النظام.

\textbf{الإجماع والإنفاذ.} يُدمج النظام الميول البشرية نحو الحقد والضغينة والرضا الانتقامي بوصفها سماتٍ حاملة للأثقال. سوء السلوك ليس محظورًا؛ بل مُسعّرًا.

% ============================================================
% 3. SYSTEM OVERVIEW
% ============================================================
\section{نظرة عامة على النظام}

\subsection{شبكة السمعة الاجتماعية}

شبكة RSN هي طبقة اتصال لامركزية من نظير إلى نظير يتبادل فيها المشاركون معلومات قابلة للتحقق حول السلوك في العالم الحقيقي. وخصائصها المميزة هي: بنية تحتية لامركزية غير قابلة للرقابة؛ ومشاركة مستعارة عبر المعرّفات اللامركزية DIDs؛ وبنية تكلفة غير متماثلة (القراءة رخيصة، والكتابة مكلفة)؛ وقابلية التحقق التشفيرية؛ و\textbf{دعم المعايير}---صيغ قابلة للقراءة آليًا للمعلومات المنتشرة عبر الشبكة، وللخدمات التي تقدّمها السلطات (تكوين الادعاء، والمصادقة، وخدمة الشهود)، ولصيغة السياسة بما في ذلك قواعد تقييم سمعة الطرف المقابل وتقدير خطر عدم تطابق السياسة (بين السلوك المُعلَن والفعلي).

\subsection{المكوّنات المعمارية}

تتألف شبكة RSN من أربعة مكوّنات أساسية (الشكل~\ref{fig:architecture}):
\begin{enumerate}[nosep]
\item \textbf{طبقة DID.} هويات المشاركين مع القواعد المُعلَنة، وبيانات السمعة الوصفية، وتوجيه الشبكة.
\item \textbf{بروتوكول الادعاء.} صيغة منظمة لنشر التأكيدات حول أحداث العالم الحقيقي.
\item \textbf{شبكة التحقق.} بروتوكول لامركزي لتعيين المحققين باستخدام التجزئة المتسقة.
\item \textbf{طبقة تجميع السمعة.} خدمات تُنتج ملخصات سمعة قابلة للقراءة البشرية.
\end{enumerate}

\begin{figure}[ht]
\centering
\begin{tikzpicture}[
    layer/.style={draw, minimum height=0.7cm, align=center, font=\scriptsize, fill=black!8},
    arr/.style={<->, thick, gray!60}
]
\node[layer, minimum width=5cm] (agg) at (0,2.8) {\textbf{التجميع}\\التفسير $\cdot$ الخطر};
\node[layer, minimum width=2.2cm] (claim) at (-1.55,1.4) {\textbf{الادعاءات}\\التكوين};
\node[layer, minimum width=2.2cm] (verif) at (1.55,1.4) {\textbf{التحقق}\\الاختيار};
\node[layer, minimum width=5cm] (did) at (0,0) {\textbf{طبقة DID}\\الهوية $\cdot$ القواعد $\cdot$ الدرجات};

\draw[arr] (agg.south west) -- (claim.north);
\draw[arr] (agg.south east) -- (verif.north);
\draw[arr] (claim.east) -- (verif.west);
\draw[arr] (claim.south) -- (did.north west);
\draw[arr] (verif.south) -- (did.north east);
\end{tikzpicture}
\caption{معمارية RSN ذات الطبقات الأربع.}
\label{fig:architecture}
\end{figure}

\subsection{أدوار المشاركين}

تتضمن معاملة التحقق ما يصل إلى ستة أدوار DID متميزة (الشكل~\ref{fig:roles}): \textbf{المُصدِر} ($DID_I$، ينشئ الادعاء ويقدّمه)، \textbf{الموضوع} ($DID_S$، المعرّف DID الذي يدور حوله الادعاء---قد يتطابق مع المُصدِر)، \textbf{السلطة} ($DID_A$، تصادق على الادعاء مقابل رسم؛ وقد تقدّم أيضًا خدمات الشهادة---\emph{اختياري})، \textbf{الشاهد} ($DID_{W_{1..n}}$، مراقب مجهول لنزاهة المحقّق عبر رموز التحدّي---\emph{اختياري})، \textbf{المُحقِّق} ($DID_V$، يُختار خوارزميًا لنشر الادعاء)، و\textbf{الشبكة RSN} (الشبكة التي تُنشَر فيها الادعاءات المُتحقَّق منها). ويمكن تفويض أي دور إلى معرّف DID آخر (القسم~7.4). وعادةً ما تجمع السلطة بين المصادقة وخدمات الشهادة، لكن الوظيفتين مستقلتان منطقيًا.

\begin{figure}[ht]
\centering
\begin{tikzpicture}[
    role/.style={draw, rounded corners, minimum width=1.1cm, minimum height=0.45cm, align=center, font=\scriptsize, fill=black!8},
    opt/.style={draw, rounded corners, minimum width=1.1cm, minimum height=0.45cm, align=center, font=\scriptsize, fill=black!8, dashed},
    arr/.style={-{Stealth[length=3pt]}, thick},
    oarr/.style={-{Stealth[length=3pt]}, thick, dashed, gray},
    lbl/.style={font=\tiny, fill=white, inner sep=1pt}
]
\node[role] (I) at (0,2.4) {\textbf{المُصدِر}};
\node[role] (S) at (-2.5,2.4) {\textbf{الموضوع}};
\node[opt] (A) at (2.5,1.2) {\textbf{السلطة}};
\node[opt] (W) at (-2.5,0) {\textbf{الشاهد}};
\node[role] (V) at (2.5,-0.6) {\textbf{المُحقِّق}};
\node[role, fill=black!5, draw=gray] (N) at (0,-1.8) {RSN};

\draw[arr] (I) -- node[lbl] {حول} (S);
\draw[oarr] (I) -- node[lbl,sloped] {يصادق} (A);
\draw[oarr] (I) -- node[lbl,sloped] {تحدٍّ} (W);
\draw[arr] (I) -- node[lbl,sloped] {طلب} (V);
\draw[oarr] (W) -- node[lbl] {يراقب} (V);
\draw[arr] (V) -- node[lbl] {ينشر} (N);
\end{tikzpicture}
\caption{أدوار DID في معاملة تحقق. المتقطع = اختياري (السلطة، الشاهد).}
\label{fig:roles}
\end{figure}

\subsection{عدم القابلية للفساد: أربع قوى}

لا تنبع مناعة RSN ضد الفساد من آلية واحدة بل من أربع قوى متزامنة. لا تكفي أيٌّ منها وحدها؛ وحده اجتماعها يجعل محاولة الفساد غير عقلانية اقتصاديًا.

\textbf{هدف غير متوقع.} يُختار المحقّق لكل ادعاء بشكل غير حتمي عبر التجزئة المتسقة (القسم~5.3). لا يستطيع المهاجم استهداف محقّق بعينه مسبقًا للرشوة، لأن هوية المحقّق دالةٌ في محتوى الادعاء والتكرارات السابقة، لا في اختيار المُصدِر.

\textbf{الرافعة السمعية---صعود بطيء، هبوط سريع.} لا يمكن اكتساب السمعة إلا تدريجيًا، عبر سلوك متسق مستدام. غير أنها يمكن أن تُفقَد كارثيًا في مصادقة احتيالية واحدة (القسم~6.2). يعني هذا اللاتماثل أن سنوات من السمعة المتراكمة يمكن أن تُدمَّر بمصادقة غير أمينة واحدة؛ فتظل الاستراتيجية المثلى هي رفض الادعاءات المشبوهة.

\textbf{الوعد العلني.} يُعلن كل حامل DID سياسته علنًا---مجموعة قواعد قابلة للقراءة آليًا يلتزم باتباعها. والتباعد بين الإعلان والسلوك الفعلي قابل للكشف ويُسعّر بوصفه مخالفة سمعة أشد وطأة من الانتهاك الأساسي نفسه (القسم~7.3). ولذلك يكون النفاق أعلى تكلفة من عدم الأمانة المباشر.

\textbf{لاتماثل تكلفة الهجوم الشبكي.} تنمو تكلفة الرقابة على الشبكة أو زعزعتها بشكل غير خطي مع حجم الشبكة وكثافتها، بينما تظل تكلفة تحمّلها ثابتة. ولكي تؤتي الرقابة ثمارها، سيتعيّن على فاعل مركزي تطبيقها على الشبكة بأكملها---مما يتطلب تدابير لا تتناسب مع أي منفعة يمكن تصوّرها (القسم~10).

% ============================================================
% 4. DECENTRALIZED IDENTITY MODEL
% ============================================================
\section{نموذج الهوية اللامركزية}

\subsection{معرّف DID بخصائص خاصة}

توسّع شبكة RSN مواصفة W3C DID Core~\cite{did2022} بما يلي: \textbf{القواعد المُعلَنة} (التزامات سلوكية قابلة للقراءة آليًا)، \textbf{بيانات السمعة الوصفية} (درجة سلوكية مجمَّعة)، و\textbf{توجيه الشبكة} (بوابة بصلية قابلة للتغيير منفصلة عن الموضع الثابت على الحلقة).

\subsection{دورة حياة الادعاء}

يمرّ الادعاء بست مراحل (الشكل~\ref{fig:lifecycle}): التكوين، والمصادقة الاختيارية من السلطة، واختيار المحقّق عبر التجزئة المتسقة، وقرار التحقق (القبول أو الرفض مع التكرار)، والنشر، وتحديث السمعة لجميع الأطراف.

\begin{figure}[ht]
\centering
\begin{tikzpicture}[
    stp/.style={draw, rounded corners=2pt, minimum width=1.3cm, minimum height=0.45cm, align=center, font=\tiny, fill=black!5},
    arr/.style={-{Stealth[length=3pt]}, thick},
    lbl/.style={font=\tiny, fill=white, inner sep=1pt}
]
\node[stp] (comp) at (0,0) {تكوين\\الادعاء};
\node[stp, dashed] (auth) at (2.0,0) {مصادقة\\السلطة};
\node[stp] (hash) at (4.0,0) {تجزئة و\\تعيين الحلقة};
\node[stp] (ver) at (2.0,-1.6) {فحص\\المحقّق};
\node[stp, double] (pub) at (0,-1.6) {مُنشَر};
\node[stp] (rej) at (4.0,-1.6) {مرفوض\\التكلفة $\uparrow$};

\draw[arr] (comp) -- (auth);
\draw[arr] (auth) -- (hash);
\draw[arr] (hash) -- (ver);
\draw[arr] (ver) -- node[lbl] {\checkmark} (pub);
\draw[arr] (ver) -- node[lbl] {$\times$} (rej);
\draw[arr] (rej) -- ++(0,0.8) -| (hash);
\end{tikzpicture}
\caption{دورة حياة الادعاء مع حلقة التكرار.}
\label{fig:lifecycle}
\end{figure}

\subsection{العلاقة بـ W3C DID Core}

نموذج الهوية في RSN هو مجموعة فوقية صحيحة لمواصفة W3C DID Core~\cite{did2022}. فجميع معرّفات DID في RSN هي معرّفات DID صالحة وفق W3C. والامتدادات الخاصة بـ RSN مُرمَّزة بوصفها خصائص لمستند DID مُعرَّفة في مواصفة طريقة DID مخصصة، متوافقة مع البنية التحتية القائمة مثل ION~\cite{buchner2021} و Ceramic~\cite{ceramic2023}.

% ============================================================
% 5. VERIFICATION AND PROPAGATION
% ============================================================
\section{التحقق من المعلومات وانتشارها}

\subsection{تكلفة إدخال المعلومات}

المبدأ الاقتصادي الأساسي لـ RSN هو اللاتماثل بين القراءة والكتابة. فقراءة بيانات السمعة تقترب من تكلفة حدّية صفرية للمشاركين الذين يشكّلون جزءًا من شبكة DID ولديهم وصول يتناسب مع سمعتهم---أما الوافدون الجدد فعليهم اكتساب وصول أوسع تدريجيًا (انظر مكافحة الاستغلال الطفيلي). أما الكتابة---المفهومة بوصفها التجسيد الدائم للسمعة في الشبكة، لا بوصفها فعلًا تقنيًا لمرة واحدة---فتتطلب استثمارًا تراكميًا قابلًا للتحقق عبر ثلاثة أبعاد: \textbf{الوقت} (سنوات من العمر أُنفقت على سلوك متسق، والتزامات مُنجَزة، وعلاقات مُنمّاة)، \textbf{الطاقة} (الجهد المستثمر في التعامل الأمين، والعناية بالشراكات، والجدارة بالثقة على المدى الطويل)، و\textbf{المال} (الاستثمار في اسم المرء نفسه---التطوير المهني، وجودة العمل، وتعويض الأضرار المتسبَّب فيها). لا يمكن شراء السمعة لحظيًا---إنها نتاج تراكم مدى الحياة يقتصر النظام على جعله مرئيًا وقابلًا للنقل عبر السياقات. والتكاليف التقنية لمرة واحدة في البروتوكول (التوقيعات التشفيرية، ورسوم المحققين) لا تُذكر مقارنة بهذا الاستثمار الأساسي.

\subsection{الرسم البياني الاجتماعي وعمق الاتصال}

RSN في جوهرها شبكة اجتماعية: يضيف كل معرّف DID جهات اتصال (معرّفات DID أخرى) تسمح بالاتصال. ولهذه الجهات جهات اتصالها الخاصة، وهكذا. ويعمل البحث الخوارزمي عن المحقّق ضمن عمق قابل للضبط $h$ من جهات الاتصال المرتبطة (مثلًا $h=3$: جهات اتصال المرء المباشرة، وجهات اتصالهم، وجهات اتصال جهات اتصالهم). لا تتطلب هذه المعمارية سلسلة كتل عالمية واحدة---بل تُحابي طبيعيًا نشوء المجتمعات/التجمّعات مع تداخلات في مجتمعات أخرى. ويجب أن يكون لكل مشارك DID \textbf{سياسة} منشورة---مجموعة قواعد قابلة للقراءة آليًا تحكم كيفية تقييم الطلبات الواردة. وتُسجَّل انتهاكات السياسة في الشبكة بوصفها آثارًا سلبية.

\subsection{الاختيار الخوارزمي للمحقّق}

يستخدم بروتوكول التحقق التجزئة المتسقة~\cite{karger1997} (الشكل~\ref{fig:verifier}). التحقق خدمة مدفوعة: يعمل المحققون مقابل رسم، مما يخلق سوقًا تدفع فيه المنافسة إلى تحسين التسعير والسرعة والموثوقية.

\textbf{الخطوة 1.} يُدمَج مستند الادعاء مع نتيجة تجزئة التكرار السابق (أو $\varnothing$ في التكرار الأول) ويُجزَّأ:
\begin{equation}
\text{pos} = f_h(\text{claim} \| \text{prev\_hash})
\label{eq:hash}
\end{equation}

يجب أن تتضمن الخوارزمية \emph{المسار الكامل} لجميع الطلبات والردود السابقة---لا مجرد تجزئة للتكرار السابق. فتُلحَق استجابة كل محقّق الكاملة (القبول، والرفض، والسعر المعروض، والسبب) بالمستند المتنامي، وهي قابلة للتحقق عبر التوقيعات التشفيرية في كل خطوة.

\textbf{الخطوة 2.} تُعيَّن التجزئة إلى $N$ موضعًا على حلقة التجزئة المتسقة، موزّعة بالتساوي (مثلًا لـ $N=4$: $+0^{\circ}, +90^{\circ}, +180^{\circ}, +270^{\circ}$). ومن كل موضع، يختار بحثٌ في اتجاه عقارب الساعة أقرب معرّف DID بوصفه مرشحًا للمحقّق (الشكل~\ref{fig:multipoint}). و$N$ معامل متغير قد يعتمد على نسبة معرّفات DID النشطة حاليًا، أو وقت اليوم، أو استجابة الشبكة التاريخية.

\textbf{الخطوة 3.} يقبل المحقّق (فيَنشر، مراهنًا بسمعته) أو يرفض (فيعود إلى الخطوة~1 بتجزئة الرفض). وعندما لا تستجيب عقدة رغم إعلانها حالة الاتصال، يجب أن يتضمن الطلب التالي جميع استجابات المرشحين المحققين الـ $N$ السابقين.

\textbf{مكافحة الاستغلال الطفيلي.} يجوز لمعرّفات DID المستهدفة أن تفرض رسومًا أو قيود وصول على الاستعلامات الواردة من معرّفات DID جديدة قليلة السمعة. وتُجبر هذه الآلية المتدرّجة (لا الثنائية) الوافدين الجدد على بناء السمعة عبر نشاط حقيقي قبل استهلاك بيانات الآخرين بحرية.

\begin{figure}[ht]
\centering
\begin{tikzpicture}[
    box/.style={draw, rounded corners=2pt, minimum width=1.3cm, minimum height=0.45cm, align=center, font=\scriptsize, fill=black!5},
    dec/.style={draw, diamond, aspect=2.2, minimum width=0.9cm, align=center, font=\scriptsize, fill=black!5},
    arr/.style={-{Stealth[length=3pt]}, thick},
    lbl/.style={font=\tiny, fill=white, inner sep=1pt}
]
\node[box] (iss) at (0,0) {المُصدِر};
\node[box] (hash) at (0,-1.1) {$f_h$(claim $\|$\\prev\_hash)};
\node[box] (ring) at (0,-2.2) {تعيين الحلقة\\عقارب الساعة};
\node[box, dashed] (spam) at (0,-3.2) {وديعة\\مكافحة البريد المزعج};
\node[dec] (check) at (0,-4.4) {فحص\\السياسة};
\node[box] (rej) at (2,-5.6) {التكرار\\التالي};
\node[box] (fee) at (0,-5.6) {الرسم النهائي =\\$f$(البيانات، السمعة، السياسة)};
\node[box, double] (pub) at (0,-6.8) {مُنشَر};

\draw[arr] (iss) -- (hash);
\draw[arr] (hash) -- (ring);
\draw[arr] (ring) -- (spam);
\draw[arr] (spam) -- (check);
\draw[arr] (check) -- node[lbl] {$\times$} (rej);
\draw[arr] (check) -- node[lbl] {\checkmark} (fee);
\draw[arr] (fee) -- (pub);
\draw[arr] (rej.east) -- ++(0.5,0) |- (hash.east);
\end{tikzpicture}
\caption{بروتوكول اختيار المحقّق. وديعة مكافحة البريد المزعج اختيارية وقد تكون قابلة للاسترداد. ولا يُدفع الرسم النهائي إلا عند القبول.}
\label{fig:verifier}
\end{figure}

يُلحق كل رفضٍ استجابة المحقّق الكاملة (بما فيها الأسباب والشروط) بمستند الادعاء، فيتوسّع محتواه. ومن المستند الموسّع، تُحسَب تجزئة جديدة بشكل غير حتمي، فتُعيَّن إلى موضع مختلف على الحلقة وتختار محقّقًا مختلفًا. وتوثيق المسار الكامل إلزامي---فلا يستطيع المُصدِر التنبؤ بالمحقّق التالي أو التأثير فيه. وإذا تجاهل محقّقٌ طلبًا رغم سياسة مُعلَنة متوافقة، سجّل ذلك الشهود (إن وُجدوا)، وجاز للمُصدِر إرفاق دليل الشهود عند النشر النهائي.

\begin{figure}[ht]
\centering
\begin{tikzpicture}[scale=0.65]
% Ring
\draw[thick] (0,0) circle (2.2cm);
% DID nodes on ring
\foreach \a in {10,55,80,120,150,195,230,260,305,340} {
  \fill[black!40] (\a:2.2) circle (1.5pt);
}
% Hash offset arrow pointing to initial position
\draw[-{Stealth[length=3pt]}, thick, black!60] (0,0) -- node[font=\tiny,above,sloped,pos=0.6] {$f_h$} (37:2.2);
\fill[black] (37:2.2) circle (2pt);
\node[font=\tiny,anchor=south west] at (37:2.35) {$h_0$};
% N=4 points offset from h0 by +0, +90, +180, +270
\foreach \offset/\l in {0/P1,90/P2,180/P3,270/P4} {
  \pgfmathsetmacro{\ang}{37+\offset}
  \node[fill=black, circle, inner sep=1.5pt] (p\offset) at (\ang:2.2) {};
  \node[font=\tiny,font=\bfseries] at (\ang:2.75) {\l};
  % clockwise search arc
  \draw[-{Stealth[length=2.5pt]}, thick, gray] (\ang:2.2) arc (\ang:\ang+30:2.2);
}
\node[font=\scriptsize, align=center] at (0,0) {حلقة\\التجزئة};
\node[font=\tiny, align=center] at (0,-3.2) {$h_0 = f_h(\text{claim})$، ثم $+0^{\circ},+90^{\circ},+180^{\circ},+270^{\circ}$\\كل نقطة $\rightarrow$ بحث باتجاه عقارب الساعة};
\end{tikzpicture}
\caption{الاختيار متعدد النقاط ($N{=}4$). التجزئة $h_0$ تحدّد الإزاحة؛ 4 نقاط موزّعة بالتساوي من $h_0$.}
\label{fig:multipoint}
\end{figure}

\subsection{بروتوكول الشهود}

يوفّر دور \textbf{الشاهد} مساءلةً مجهولة لنزاهة المحقّق عبر بروتوكول رمز تحدٍّ تشفيري:

\textbf{الخطوة W1.} يوقّع الطالب طلب التحقق ويرسله إلى $N_w$ شاهدًا---جهات اتصال من شبكة الطالب الاجتماعية، أو مقدّمي خدمات شهادة متخصصين يعملون مقابل رسم.

\textbf{الخطوة W2.} يوقّع كل شاهد $w_i$ على الطلب المُوقَّع بأكمله، ويحتفظ بالتوقيع الأصلي $\sigma_i$، ويحسب رمز تحدٍّ $c_i = h(\sigma_i)$. ويُلحَق رمز التحدّي بالطلب.

\textbf{الخطوة W3.} يُرسَل الطلب، وهو يحمل الآن $N_w$ رمز تحدٍّ، إلى المحقّق. يشاهد المحقّق الرموز لكنه \emph{لا يستطيع تحديد} من هم الشهود أو ما إذا كانت الرموز تقابل توقيعات حقيقية.

\textbf{الخطوة W4 (محقّق أمين).} إذا استجاب المحقّق وفق سياسته المُعلَنة، تظل رموز التحدّي معتمة. ولا يُكشف الشهود أبدًا. ولا تُنشر بيانات إضافية.

\textbf{الخطوة W5 (انتهاك السياسة).} إذا انتهك المحقّق سياسته (تجاهل الطلب، أو استجاب بشكل غير متسق)، فإن الطالب يحتفظ بتوقيعات الشهود الأصلية $\sigma_i$. ويمكن نشرها بوصفها بيانات مصاحبة: يستطيع أي شخص التحقق من $c_i = h(\sigma_i)$، مثبتًا حضور الشاهد. ويستطيع الطالب النشر بشكل مستقل---فقد وقّع المحقّق سلفًا على رموز التحدّي في استجابته.

\textbf{خاصية الخديعة.} يجوز للطالب استبدال قيم عشوائية ببعض رموز التحدّي أو كلها. لا يستطيع المحقّق تمييز الرموز الحقيقية (المدعومة بسلطات عالية السمعة) عن الضجيج. هذا \emph{اللايقين} هو آلية الإنفاذ: فكل طلب يحمل خطر أن تكون سلطة محترمة تراقب متخفّية. وتكلفة الحفاظ على هذا الضغط تقارب الصفر (توليد أرقام عشوائية)، بينما التكلفة المحتملة لعدم الأمانة على المحقّق كارثية. فيُحفَّز السلوك الأمين حتى في غياب شهود فعليين.

\textbf{السلطة بوصفها شاهدًا.} يجوز لسلطة تصادق على ادعاء أن تجمع خدمات الشهادة: «أصادق على ادعائك وأعمل شاهدًا متخفّيًا أثناء التحقق». هذا نموذج عمل طبيعي---فالسلطة تملك السياق أصلًا. غير أن الدورين مستقلان منطقيًا.

\subsection{مبدأ التطرف المكلف}

تتراكم ثلاث قنوات تكلفة في آنٍ واحد (الشكل~\ref{fig:radicalism}):

\textbf{القناة 1: تكلفة التكرار.} يفرض كل رفضٍ تكرارًا جديدًا يستهلك الوقت والموارد. نموّ خطي.

\textbf{القناة 2: تضخّم المستند.} يضيف كل تكرارٍ استجابة المحقّق الكاملة إلى مستند الادعاء. النموّ خطي، لكن مع إزاحة أساسية: فالحمولة الأولية (محتوى الادعاء، ومصادقة السلطة، والتوقيعات التشفيرية) لها بالفعل حجم غير تافه $B_0$. الحجم الإجمالي بعد $k$ تكرارًا: $S(k) = B_0 + k \cdot \Delta$، حيث $\Delta$ متوسط حجم الاستجابة.

\textbf{القناة 3: علاوة مخاطر المحقّق.} الخطر ذاتي. يقيّم المحقّق ما إذا كانت السياسات المُعلَنة لمعرّف DID الطالب تتعارض مع مصالحه التجارية أو الاجتماعية أو السياسية. وقد تتصاعد العلاوة أسّيًا مع البُعد المُدرَك عن «مثال» المحقّق: $P(d) = \alpha \cdot e^{\beta d}$، حيث $d$ مقياس المسافة الذاتي للمحقّق. وبعد حدٍّ شخصي محظور، قد يرفض المحقّق كليًا.

\begin{figure}[ht]
\centering
\begin{tikzpicture}
\begin{axis}[
    width=\columnwidth,
    height=4.5cm,
    xlabel={\footnotesize التطرف},
    ylabel={\footnotesize التكلفة النسبية (لوغاريتمي)},
    ymode=log,
    xmin=0, xmax=100,
    ymin=1, ymax=10000,
    grid=major,
    grid style={gray!20},
    legend style={at={(0.03,0.97)},anchor=north west,font=\tiny,draw=gray!50},
    tick label style={font=\tiny},
    label style={font=\footnotesize},
]
\addplot[black, thick, domain=0:100, samples=50] {1 + 0.3*x};
\addlegendentry{تكلفة التكرار (خطية)}
\addplot[black, thick, dashed, domain=0:100, samples=50] {1 + 0.15*x};
\addlegendentry{تضخّم المستند (خطي)}
\addplot[black, thick, dotted, line width=1.2pt, domain=0:100, samples=100] {exp(0.08*x)};
\addlegendentry{علاوة المخاطر (أسّية)}
\end{axis}
\end{tikzpicture}
\caption{تصاعد التكلفة عبر ثلاث قنوات. علاوة المخاطر تهيمن عند التطرف العالي.}
\label{fig:radicalism}
\end{figure}

والأهم أن هذا ليس رقابة. فلا يُحظَر أي ادعاء. النظام يُسعّر المخاطر (الجدول~\ref{tab:costs}).

\begin{table}[ht]
\centering
\caption{مقارنة التكلفة عبر أنواع الادعاءات.}
\label{tab:costs}
\footnotesize
\begin{tabular}{@{}lcccc@{}}
\toprule
\textbf{النوع} & \textbf{التكرار} & \textbf{المستند} & \textbf{الرسم} & \textbf{الإجمالي} \\
\midrule
موثوق & $k_{\min}$ & $B_0$ & $P_{\min}$ & منخفض \\
مثير للجدل & $k_{\text{med}}$ & $B_0{+}k\Delta$ & $\alpha e^{\beta d}$ & متوسط \\
جذري & $k_{\max}$ & $\gg B_0$ & $\gg P_{\min}$ & مرتفع جدًا \\
احتيالي & $\to\infty$ & --- & --- & غير قابل للنشر \\
\bottomrule
\end{tabular}
\end{table}

% ============================================================
% 6. REPUTATION AND RISK
% ============================================================
\section{السمعة وتقييم المخاطر}

\subsection{نموذج تراكم السمعة}

تُظهر السمعة ثلاث خصائص: \textbf{تراكم بطيء} (نموّ تدريجي عبر سلوك متسق)، \textbf{تدمير سريع} (احتيال واحد قد يخفض الدرجة كارثيًا)، و\textbf{تقييم فردي} (لكل طرف مستهلك أن يطبّق دالة تجميعه الخاصة).

\subsection{السلطات ذات السمعة}

تراجع السلطة ذات السمعة الادعاءات، وإن اقتنعت، تُشارك في توقيعها---مراهنةً بسمعتها الخاصة (الشكل~\ref{fig:authority}). واللاتماثل مقصود بالتصميم: فاكتساب السمعة بطيء (زيادة تدريجية $+\delta$ لكل مصادقة أمينة)، بينما فقدانها كارثي ($-\Delta \gg \delta$ لكل مصادقة كاذبة؛ على سبيل الإيضاح، مثلًا $+2\%$ مقابل $-70\%$). والاستراتيجية المثلى هي دائمًا رفض الادعاءات المشبوهة. أما القيم المحددة لـ $\delta$ و $\Delta$ فهي معاملات النظام.

\begin{figure}[ht]
\centering
\begin{tikzpicture}[
    box/.style={draw, rounded corners=2pt, minimum width=2cm, minimum height=0.6cm, align=center, font=\scriptsize, fill=black!5},
    arr/.style={-{Stealth[length=4pt]}, thick}
]
\node[box] (exam) at (0,0) {السلطة تفحص\\السمعة: $R$};
\node[box] (true) at (-2,-1.3) {صدق: $R{\to}R{+}\delta$\\عملاء أكثر};
\node[box] (false) at (2,-1.3) {كذب: $R{\to}R{-}\Delta$\\العملاء يغادرون};
\node[box] (insight) at (0,-2.6) {الكسب: بطيء ($+\delta$)\\الخسارة: فورية ($-\Delta$)};

\draw[arr] (exam) -- node[left,font=\tiny] {صدق} (true);
\draw[arr] (exam) -- node[right,font=\tiny] {كذب} (false);
\draw[arr] (true) -- (insight);
\draw[arr] (false) -- (insight);
\end{tikzpicture}
\caption{السلطة ذات السمعة---حوافز غير متماثلة.}
\label{fig:authority}
\end{figure}

\subsection{السمعة متعددة الأبعاد}

السمعة ليست كمية قياسية بل متجهًا عبر أبعاد متعددة: الموثوقية المالية، والنزاهة الشخصية، والكفاءة المهنية، والمشاركة المدنية، وغيرها (الشكل~\ref{fig:reputation-radar}). ويُسقط مقدّمو التقييم المختلفون هذا المتجه بشكل مختلف تبعًا للسياق---فالمُقرِض يعطي الأولوية للموثوقية المالية، وصاحب العمل للكفاءة المهنية، والجار للسلوك المدني. ولا توجد درجة سمعة «صحيحة» واحدة؛ بل مجرد منظورات.

\begin{figure}[ht]
\centering
\begin{tikzpicture}[scale=0.55]
\foreach \a/\l in {90/مالية,162/نزاهة,234/مهنية,306/مدنية,18/اجتماعية} {
  \draw[gray!40] (0,0) -- (\a:2.5);
  \node[font=\tiny, align=center] at (\a:3.1) {\l};
  \foreach \r in {0.8,1.6,2.4} {
    \fill[black!15] (\a:\r) circle (0.5pt);
  }
}
\draw[thick, fill=black!10] (90:2.0) -- (162:1.2) -- (234:1.8) -- (306:0.9) -- (18:1.5) -- cycle;
\draw[thick, dashed] (90:1.4) -- (162:2.1) -- (234:0.8) -- (306:2.0) -- (18:1.1) -- cycle;
\node[font=\tiny] at (0,-3.3) {مصمت: DID A \quad متقطع: DID B};
\end{tikzpicture}
\caption{متجهات السمعة متعددة الأبعاد. تُظهر معرّفات DID المختلفة ملامح مختلفة عبر الأبعاد.}
\label{fig:reputation-radar}
\end{figure}

\subsection{تقييم المخاطر المُدمقرَط}

تُدمقرِط شبكة RSN تقييم المخاطر المنهجي---وهو قدرة تتركّز حاليًا في المؤسسات المالية لكنها تنطبق على ما هو أبعد بكثير من الصيرفة. فتقييم التكتلات الصناعية لموثوقية المورّدين، وتسعير شركات التأمين للمخاطر السلوكية، والعناية الواجبة في عمليات الاندماج والاستحواذ، وتقدير عمر المعدات في التصنيع---أينما تطلّب خطر الطرف المقابل تقييمًا، توفّر شبكة RSN بديلًا لامركزيًا مدفوعًا بالسوق. وتُترجم سوق خدمات التفسير بيانات السمعة الخام متعددة الأبعاد إلى ملخصات قابلة للتنفيذ، مما يجعل تقييم الطرف المقابل متاحًا لأي مشارك بتكلفة حدّية.

% ============================================================
% 7. CONSENSUS AND DISPUTE RESOLUTION
% ============================================================
\section{الإجماع وحل النزاعات}

\subsection{نموذج العدالة اللامركزية}

تعيد شبكة RSN صياغة العدالة بوصفها مسألة تفاوض ثنائي. فالتهديد بضرر سمعي دائم وقابل للتحقق رادعٌ أكثر فعالية من الإجراءات القانونية التي لا يقدر الطرف المتضرر على تحمّلها.

\subsection{الإجماع الناشئ}

يحتفظ كل مشارك بعتبة رضا شخصية. وينشأ الإجماع الكلي للشبكة بوصفه المتوسط الإحصائي لآلاف المفاوضات الثنائية (الشكل~\ref{fig:consensus}). فالاستجابة الأعلى بكثير من الإجماع (همجية) مكلفة النشر. والاستجابة القريبة من الإجماع (متناسبة) رخيصة. والإجماع ليس قاعدة---إنه إشارة سعرية.

\begin{figure}[ht]
\centering
\begin{tikzpicture}[
    person/.style={draw, rounded corners=2pt, minimum width=1.5cm, minimum height=0.5cm, align=center, font=\scriptsize, fill=black!5},
    arr/.style={-{Stealth[length=4pt]}, thick}
]
\node[person] (a) at (-2,0) {صارم\\(عالٍ)};
\node[person] (b) at (0,0) {عملي\\(متوسط)};
\node[person] (c) at (2,0) {متسامح\\(منخفض)};
\node[person, fill=black!10, minimum width=3cm] (cons) at (0,-1.2) {إجماع الشبكة\\= متوسط إحصائي};
\node[person] (exp) at (-2,-2.4) {همجي\\مكلف};
\node[person, double] (prop) at (0,-2.4) {متناسب\\رخيص};
\node[person] (neg) at (2,-2.4) {متهاون\\مكلف};

\draw[arr] (a) -- (cons);
\draw[arr] (b) -- (cons);
\draw[arr] (c) -- (cons);
\draw[arr] (cons) -- (exp);
\draw[arr] (cons) -- (prop);
\draw[arr] (cons) -- (neg);
\end{tikzpicture}
\caption{الإجماع الناشئ من عتبات الرضا الفردية.}
\label{fig:consensus}
\end{figure}

\subsection{معايرة العقاب}

يُعلن كل حامل DID علنًا استجاباته لفئات محددة من سوء السلوك. والإعلانات القاسية أو المتساهلة بشكل غير متناسب تجذب إشارات سمعة سلبية. أما الإعلانات المتناسبة فتُقبَل دون احتكاك---نظام عقاب ذاتي المعايرة.

وتخضع إعلانات الإجماع ذاتها لكشف النفاق: فالالتزام إعلانيًا بموقف إجماعي مع التصرف بشكل غير متسق يشكّل مخالفة سمعة أشد وطأة من الانحراف الأساسي، لأنه يبرهن على خداع متعمّد.

\subsection{التفويض}

يجوز تفويض جميع الأنشطة ضمن معرّف DID---باستثناء واحد---إلى معرّف DID آخر. \textbf{دور الموضوع---المعرّف DID الذي يدور الادعاء حول سلوكه---لا يمكن تفويضه؛ فهو مرتبط ارتباطًا غير قابل للنقل بحامل الهوية الذي يشير إليه الادعاء.} أما الأدوار النشطة الأخرى---المُصدِر، والسلطة، والشاهد، والمُحقِّق---فيمكن نقلها إلى معرّف DID مفوَّض مُعيَّن، سواء للتحقق، أو تقديم الادعاء، أو المشاركة في الإجماع، أو حل النزاعات. والتفويض قابل للإلغاء في أي وقت. ويُتيح هذا التخصص (مثل خدمات التحقق المهنية) مع احتفاظ الحامل بالسيطرة والمسؤولية النهائيتين. والتفويض ينطبق على النظام بأكمله وهو جوهري لقابلية التوسّع.

\subsection{من الأخلاق الفردية إلى العقد الاجتماعي الناشئ}

تمثّل السياسة المُعلَنة لكل معرّف DID تجسيدًا رسميًا للأخلاق الفردية: ما يعتبره الحامل مقبولًا، وكيف يستجيب لسلوكيات محددة، وما يطلبه من الأطراف المقابلة. وهذه السياسات ليست مفروضة من مصمّم النظام---بل يختارها كل مشارك ويعبّر عنها بصيغة قابلة للقراءة آليًا.

يُتيح هذا التجسيد الرسمي نشوءًا ثلاثي المراحل. أولًا، تُرمَّز الأخلاق الفردية بوصفها \textbf{سياسة}---مجموعة قواعد وعتبات ودوال استجابة مُعلَنة يلتزم بها المعرّف DID. ثانيًا، يُنتج مجموع كل السياسات عبر الشبكة \textbf{أخلاقًا ناشئة}---أعرافًا قابلة للرصد إحصائيًا لم يصممها أي مشارك بمفرده لكنها تنشأ من تفاعل آلاف المواقف الأخلاقية الفردية~\cite{durkheim1893}. ثالثًا، تشكّل هذه الأخلاق الناشئة، المعبَّر عنها بوصفها أعرافًا سلوكية مشتركة مُنفَّذة عبر إشارات سعرية ثنائية، \textbf{عقدًا اجتماعيًا قابلًا للقياس}---لا بناءً نظريًا لم يوقّعه أحد~\cite{rawls1971}، بل مجموعة قواعد قابلة للرصد تجريبيًا مدعومة بسلوك فعلي.

تعكس هذه العملية نتائج نظرية الأسس الأخلاقية~\cite{haidt2012}: فالأخلاق البشرية حدسية، وتعددية، ومتغيّرة ثقافيًا. لا تتطلب شبكة RSN إجماعًا أخلاقيًا مدخلًا؛ بل تُنتج إجماعًا سلوكيًا مخرجًا. والعقد الاجتماعي الناتج ليس ثابتًا---بل يتطور مع انضمام المشاركين ومغادرتهم وتعديلهم لسياساتهم استجابةً لتغذية الشبكة الراجعة. وخلافًا لنظرية العقد الاجتماعي الكلاسيكية، هذا العقد قابل للإلغاء، وللرصد، وللإنفاذ دون سيادة.

% ============================================================
% 8. ECONOMIC MODEL
% ============================================================
\section{النموذج الاقتصادي}

وصفت الأقسام السابقة أداةً---آليات التحقق في RSN، وبنية التكلفة، والإجماع الناشئ. ويصف هذا القسم منهجيةً لتطبيق هذه الأداة على الانتقال المالي والحوكمي. الأداة عامة الغرض؛ والمنهجية أدناه تطبيق ممكن واحد.

\subsection{بنية الحوافز}

السمعة بوصفها رأس مال تخلق حوافز مباشرة للسلوك الأمين. في التشغيل العادي، يبني المشاركون السمعة طبيعيًا عبر المعاملات اليومية والالتزامات المُنجَزة. والإنفاق الأساسي يكون على الادعاءات \emph{السلبية}---تسجيل الضرر الذي يسبّبه الآخرون. ويحفّز هذا اللاتماثل السلوك المفيد الموثوق: فالأمانة لا تكلّف شيئًا إضافيًا، بينما يخلق الإيذاء أثرًا موثَّقًا سيدفع الآخرون للحفاظ عليه. أما النفاق---إعلان القواعد دون اتباعها---فهو أكثر أنماط الفشل تكلفة، لأنه يبرهن على خداع متعمّد.

\subsection{النظام المالي الموازي}

تُتيح ثلاثة مكوّنات ضغطًا تنافسيًا: (1)~\textbf{الضريبة البسيطة}---معدل نسبي واحد بلا استثناءات، أقل هامشيًا في البداية من المعدل الفعلي القائم؛ (2)~\textbf{التسجيل الإلكتروني للإنفاق (ESR)}---بقلب نموذج EET التشيكي كي يراقب المواطنون نفقات الدولة؛ (3)~\textbf{تخصيص الضرائب الموجَّه من المواطنين}---بدءًا ببضع نسب مئوية في السنة الأولى، وبلوغًا، بعد عقد، أي مستوى يتراوح من العشرات المنخفضة إلى الهجرة الكاملة نحو نظام موجَّه من المواطنين، تبعًا لمعدل التبنّي. ويعتمد الوتيرة الفعلية على سرعة نشوء بدائل السوق الحرة لخدمات الدولة، وعلى استعداد الدولة للتعاون مع الانتقال؛ ولا يلزم أن تكون دالة الوقت خطية، ولا يوجد سبب لوضع حد أعلى لمستوى التبنّي الذي قد يبلغه في نهاية المطاف.

% ============================================================
% 9. MIGRATION PATH
% ============================================================
\section{مسار الانتقال}

يمضي الانتقال عبر ثلاث مراحل (الشكل~\ref{fig:migration}): \textbf{المرحلة~1: التراكب} (RSN بوصفها طبقة معلوماتية، الدولة هي المزوّد الوحيد)، \textbf{المرحلة~2: المنافسة} (RSN تقدّم بدائل وظيفية، استخدام مزدوج للنظامين)، و\textbf{المرحلة~3: الانتقال} (RSN تتفوق على مكافئات الدولة، هجرة طوعية). والانتقال قابل للعكس في كل مرحلة.

\begin{figure}[ht]
\centering
\begin{tikzpicture}[
    phase/.style={draw, rounded corners=3pt, minimum width=1.8cm, minimum height=0.8cm, align=center, font=\scriptsize, fill=black!5},
    arr/.style={-{Stealth[length=5pt]}, thick}
]
\node[phase] (p1) at (0,0) {\textbf{المرحلة 1}\\التراكب};
\node[phase] (p2) at (2.2,0) {\textbf{المرحلة 2}\\المنافسة};
\node[phase] (p3) at (4.4,0) {\textbf{المرحلة 3}\\الانتقال};

\draw[arr] (p1) -- (p2);
\draw[arr] (p2) -- (p3);
\draw[arr, dashed, gray] (p3.south) -- ++(0,-0.6) -| node[below,font=\tiny,pos=0.25] {تراجع} (p2.south);
\end{tikzpicture}
\caption{انتقال ثلاثي المراحل---قابل للعكس في كل مرحلة.}
\label{fig:migration}
\end{figure}

آلية الضغط الأساسية مالية: فعندما يبلغ دافعو الضرائب المشروطون «أقلية ذات دلالة اقتصادية $X$»---مجموعة كبيرة بما يكفي لأن يُلحق قمعها ضررًا اقتصاديًا غير تافه ويصبح العصيان المدني تهديدًا ذا مصداقية---تدخل الدولة في التفاوض. وللإيضاح، نستخدم $X \approx 15\%$ من الناتج المحلي الإجمالي، لكن العتبة الفعلية تعتمد على الاقتصاد المحدد.

% ============================================================
% 10. SECURITY ANALYSIS
% ============================================================
\section{التحليل الأمني}

ننظر في خمس فئات من الخصوم: الفاعلون السيئون الأفراد، والجماعات المنظمة، والخصوم من الشركات، والخصوم على مستوى الدولة، والخصوم على مستوى البروتوكول.

\textbf{مقاومة سيبيل}~\cite{douceur2002}. يتطلب بناء السمعة تفاعلًا مستدامًا قابلًا للتحقق. وتتناسب تكلفة هجوم سيبيل الناجح خطيًا مع الهويات المزيفة وتربيعيًا مع مستوى السمعة المستهدف. وتشغيل معرّفات DID متعددة بالتوازي ممكن لكنه مكلف بالتصميم---فكل هوية يجب أن تراكم سمعتها بشكل مستقل عبر نشاط حقيقي. وفي المجتمعات الحرة تردع هذه التكلفة الإساءة العابرة؛ أما في الديكتاتوريات فتصبح معرّفات DID المتوازية أداة بقاء تُمكّن من المقاومة المُجزّأة، والتنسيق السري، والتنقل الأكثر أمانًا في السوق السوداء.

\textbf{الدفاع ضد التواطؤ.} أنماط المصادقة المتبادلة بين الجماعات المغلقة قابلة للكشف عبر تحليل الرسم البياني الاجتماعي. وتُخصم دوال تجميع السمعة الادعاءات الواردة من مصادر متجمّعة بإحكام.

\textbf{الخصم على مستوى الدولة.} يضمن التوجيه البصلي، وغياب البنية التحتية المركزية، وتوزيع دور المحقّق عبر قاعدة المشاركين، أن القمع يتطلب تدابير لا تتناسب مع أي منفعة.

\textbf{الخصوصية مقابل المساءلة.} الاستعارة---أرضية وسطى بين إخفاء الهوية والإفصاح عنها---تسمح لمعرّفات DID بمراكمة سمعة ذات معنى دون كشف هوية الحامل في العالم الحقيقي.

% ============================================================
% 11. PRIVACY
% ============================================================
\section{اعتبارات الخصوصية}

يقوم نموذج الخصوصية في RSN على الاستعارة. يتحكم الحامل في الإفصاح عن الهوية: أدنى (اسم مستعار خالص)، أو انتقائي (ربط بيانات اعتماد محددة)، أو كامل (ربط الهوية القانونية). والادعاءات المنشورة دائمة---يدعم النظام التصحيح السياقي لا المحو. ويجوز للحامل التخلي عن معرّف DID مخترَق والبدء من جديد، متنازلًا عن السمعة المتراكمة.

% ============================================================
% 12. RELATED WORK
% ============================================================
\section{الأعمال ذات الصلة}

توفّر مواصفة W3C DID Core~\cite{did2022} وشبكة Ceramic~\cite{ceramic2023} أساس الهوية. وأظهر EigenTrust~\cite{kamvar2003} تجميع السمعة الموزّع دون سلطة مركزية. وقدّمت SBTs~\cite{weyl2022} رموزًا اجتماعية غير قابلة للتحويل. وتختلف شبكة RSN في تشديدها على التكلفة الاقتصادية بوصفها مرشّح جودة وفي آليتها لتسعير التطرف.

وأظهرت DAOs~\cite{buterin2014} والتصويت التربيعي~\cite{buterin2019} جدوى الحوكمة اللامركزية. وتستمد شبكة RSN الإجماع من مفاوضات سعرية ثنائية لا من التصويت.

يستند المقترح إلى النظرية الأناركية الرأسمالية~\cite{rothbard1973,hoppe2001} في توفير الخدمات الخاصة لكنه يفارقها في وجهين حاسمين: فهو يُصمَّم للضغينة والدوافع الانتقامية بدلًا من افتراض العقلانية، ويقترح انتقالًا تدريجيًا بدلًا من الثورة. ولمفهوم العقود الاجتماعية الناشئة سوابق في نظرية هايك عن النظام العفوي~\cite{hayek1973}.

\textbf{الأناركية الأغورية.} تتشارك شبكة RSN أرضية مشتركة كبيرة مع الاقتصاد المضاد الأغوري~\cite{konkin2006}---لا سيما التشديد على بناء مؤسسات موازية والتبادل الطوعي. غير أن الأغورية تميل إلى افتراض المصلحة الذاتية العقلانية بوصفها آلية تنسيق كافية وكثيرًا ما تفتقر إلى مسار انتقال ملموس من هياكل الدولة القائمة. أما شبكة RSN فتُدمج صراحةً الميول السلوكية غير العقلانية وتوفّر آلية ضغط مالي للانتقال التدريجي.

\textbf{الجدارة الاستحقاقية.} قد تمثّل شبكة RSN أول وصف وظيفي لكيفية نشوء الجدارة الاستحقاقية~\cite{young1958} بوصفها خاصية نظامية لا شعارًا. فالأنظمة الجدارية التقليدية تفشل لأنها تتطلب سلطة مركزية لتعريف «الجدارة» وإنفاذها. وفي شبكة RSN، تنشأ الجدارة من التقييمات الثنائية: فمن يقدّمون قيمة يراكمون السمعة؛ ولا تقرّر أي لجنة من له جدارة.

\textbf{الملكية بوصفها امتيازًا.} تفارق شبكة RSN جذريًا المعالجات الأناركية الرأسمالية والنمساوية لحقوق الملكية بوصفها طبيعية ومطلقة. ففي إطار RSN، الملكية \emph{امتياز}---اعتراف اجتماعي يمكن، في الحالات القصوى، أن يسحبه المجتمع عندما ينتهك مشارك انتهاكًا جوهريًا الأعراف المشتركة (الخيانة، ورفض الدفاع عن المجتمع في أزمة وجودية، والاستغلال المنهجي). وهذا ليس مصادرة من الدولة بل سحب للاعتراف الاجتماعي من شبكة العلاقات الثنائية.

% ============================================================
% 13. LIMITATIONS
% ============================================================
\section{النقاش والقيود}

\textbf{الانطلاق البارد.} تعتمد قيمة شبكة RSN على آثار الشبكة؛ فيواجه المتبنّون الأوائل قيمة محدودة إلى أن تبلغ المشاركة عتبة معينة. غير أنه حتى من المراحل المبكرة، تخدم شبكة DID بوصفها مصدر معلومات تكميلي مفيد. ويُتوقَّع أن يتطور التقييم المبكر للسمعة وتقييم السلطات تدريجيًا مع تزايد التطور.

\textbf{مكافحة الاستغلال الطفيلي والتحكم في الوصول.} يجوز لمعرّفات DID المستهدفة فرض رسوم متدرّجة وقيود وصول على الاستعلامات من معرّفات DID منخفضة السمعة. وهذا ليس ثنائيًا بل مقياس متصل: فكلما قلّت سمعة معرّف DID المستعلِم، قلّت المعلومات التي يتلقاها، وطالت مدة انتظاره، وزاد ما يدفعه. وتحفّز هذه الآلية المشاركة الحقيقية على الاستهلاك السلبي.

\textbf{عدم المساواة في الوصول.} قد تُرشّح تكلفة نشر الادعاءات المظالم المشروعة الصادرة عن المشاركين المحرومين اقتصاديًا. التخفيف: كل مشارك يخدم في الوقت نفسه بوصفه محقّقًا محتملًا (أو يفوّض هذا الدور) ويكسب رسوم التحقق. وعلى أفق زمني كافٍ، ينبغي أن يقترب مشارك غير متطرف من الحياد المالي---إذ يقابل دخل التحقق تقريبًا تكاليف النشر. وهذا توقع إحصائي لا ضمانة.

\textbf{عدم المساواة في السمعة.} يراكم المتبنّون الأوائل مزايا قد تخلق حواجز أمام المتأخرين. غير أن تقييم السمعة يقوم به مقدّمو طرف ثالث قد يختارون تخفيف ميزة المتحرّك الأول عبر خوارزميات تجميعهم. وكونك الأول بسمعة جيدة ميزة اقتصادية مقارِنة مشروعة---وذلك بالتصميم.

\textbf{أسئلة مفتوحة.} تظل دوال التكلفة المثلى، ومعايير التجميع، وحوكمة البروتوكول، والتفاعل مع بيانات السمعة الاصطناعية المولَّدة بالذكاء الاصطناعي، مجالات لعمل مستقبلي.

لا تدّعي هذه الورقة أن شبكة RSN ستُنتج نتائج مثلى في كل الظروف. بل تدّعي أن شبكة RSN تمثّل بديلًا \emph{ممكنًا}---قابلًا للتنفيذ تقنيًا، ومستدامًا ذاتيًا اقتصاديًا، ومتوافقًا اجتماعيًا مع الميول السلوكية البشرية كما هي فعلًا.

% ============================================================
% 14. CONCLUSION
% ============================================================
\section{الخاتمة}

قدّمت هذه الورقة شبكة السمعة الاجتماعية، وهي بروتوكول لامركزي يُتيح تبادل سمعة مستعار وقابل للتحقق. ويضمن مبدأ التطرف المكلف تسعير الادعاءات الاحتيالية خارج المتناول دون رقابة. ويُتيح مسار الانتقال المقترح تحوّلًا تدريجيًا قابلًا للعكس من المؤسسات القائمة. ويُدمج التصميم الطبيعة البشرية كما هي---بما في ذلك الحقد، والضغينة، والرغبة في الرضا الانتقامي---بدلًا من اشتراط تحوّل أيديولوجي. والنتيجة ليست يوتوبيا بل نظامًا تكون فيه العدالة متاحة، والسمعة مكتسَبة، وتكلفة سوء السلوك يتحمّلها الطرف الذي تسبّب فيه.

% ============================================================
% REFERENCES
% ============================================================
\bibliography{references}

\end{document}
